\chapter{Installing Blendity and the first start}
\label{ch:install}

Blendity is free. You don't need Blender or anything else: the installer brings everything Blendity uses. Version
0.13.0 is a \textbf{pre-release}: it works, but it is still being refined, so expect rough edges and save your work
often.

\section{Download}
Open the Releases page, \url{https://github.com/EmilianoManaloIV/blendity/releases}, or click the link for your
computer in the project's README: it fetches the one file you need, always the newest version.

\begin{center}
\begin{tabular}{ll}
\toprule
Your computer & The file \\
\midrule
Windows 10 or 11 & \code{Blendity-Setup-Windows.exe} \\
A Mac with an Apple chip (M1 or later) & \code{Blendity-macos.dmg} \\
Linux (Ubuntu 22.04 or newer, Fedora, Mint, \ldots) & \code{Blendity-linux-x86\_64.AppImage} \\
\bottomrule
\end{tabular}
\end{center}

Not sure which Mac you have? Apple menu $>$ About This Mac: ``Chip: Apple M\ldots'' means an Apple chip. Macs with an
Intel processor aren't supported yet. Every release also carries \code{SHA256SUMS.txt} (Blendity's updater checks
downloads against it), this guide and the architecture book as PDFs, and ``Source code'' archives that are for
developers.

\section{Install}
\textbf{Windows.} Double-click \code{Blendity-Setup-Windows.exe}. If a blue box says ``Windows protected your PC'', click
\ui{More info}, then \ui{Run anyway} (Blendity isn't signed with a paid certificate yet). Pick \ui{for me only} if asked
(no administrator needed), accept the licence, keep \ui{Blendity with Blender's libraries (recommended)}, tick a
desktop shortcut if you want one, \ui{Install}, \ui{Finish}. Next time, open it from the Start menu. To uninstall:
Settings $>$ Apps $>$ Installed apps $>$ Blendity $>$ Uninstall; your projects are kept.

\textbf{Mac.} Double-click the \code{.dmg}, drag \ui{Blendity} onto \ui{Applications}, open it from Applications. The
first time macOS says it can't check the app: on macOS 15 or newer click \ui{Done}, then System Settings $>$ Privacy \&
Security $>$ \ui{Open Anyway}; on macOS 14 or older right-click Blendity $>$ \ui{Open} $>$ \ui{Open}. If macOS says
Blendity ``is damaged'', run \code{xattr -dr com.apple.quarantine /Applications/Blendity.app} in Terminal once. To
uninstall, drag it to the Bin.

\textbf{Linux.} Right-click the \code{.AppImage} $>$ Properties $>$ Permissions $>$ \ui{Allow executing file as program},
then double-click it. If nothing happens, download the plain folder (\code{blendity-linux-x86\_64.tar.gz} from the
``libraries'' release), extract it and double-click \ui{Blendity} inside. To uninstall, delete the file or the folder.

\section{The first start}
\begin{itemize}
  \item The welcome asks \textbf{where you come from} (Unity, Unreal, Blender, Photoshop \& Substance) and sets the keys and
    the editor feel together; Help $>$ Welcome asks again, and chapter~\ref{ch:comingfrom} explains each choice.
  \item Your work is saved in a \textbf{project folder}: \code{BlendityProject} in your home folder unless you open
    another. A project is just a folder with \code{Assets} inside, like Unity's; copy it to another machine and open a
    scene in it. The program's folder belongs to the installer; your scenes, materials and renders never live there.
  \item Press \key{F1} for the \textbf{Learn tab}: short interactive lessons for each part of Blendity, with Try-it
    buttons that drive the editor while you watch (chapter~\ref{ch:learn}).
  \item \textbf{Updates:} once a day, when online, Blendity asks GitHub whether a newer build exists and shows what it
    brings with \ui{Install} / \ui{Not now} / \ui{Skip this one}; installs are checked against the release's SHA-256
    list. Help $>$ Version History goes back to an older version if a new one gives you trouble. Preferences $>$ Updates
    turns the check off or limits it to stable releases.
  \item \textbf{With or without Blender's libraries:} an install holds two programs, Blendity (its own code only, always
    starts) and Blendity-full (with Blender's libraries: faster ray tracing, denoising, OpenColorIO's views, Jolt
    physics, video export, Booleans, Decimate, EXR). Blendity starts first and hands over to the full one when its
    libraries load and you want them; offline, or with a library missing, it runs on its own code and says so.
    Preferences $>$ Libraries is your choice; \code{--no-libs} starts without them once; when they are missing and you
    are online, Blendity offers to download them for this build.
  \item \textbf{On a small laptop} (the HP Envy 13-inch 2-in-1 class: 4 cores, 8~GB, Iris Xe) Blendity starts on the
    \ui{Low} quality tier and says so, with \ui{High} one click away (chapter~\ref{ch:trouble}).
\end{itemize}

\section{Starting from the command line}
\begin{lstlisting}[language=console]
Blendity --layout Learning            # start with the Learn tab beside the Scene view
Blendity --scene Assets/Scenes/SampleScene.scene
Blendity --scale 1.5                  # force the UI scale (the default follows the OS DPI)
Blendity --headless-screenshot out.png --cmd "select Cube" --cmd "shading rendered" --frames 30
\end{lstlisting}
The last form runs with no window: it executes console commands and writes a picture, which is how the tests, the
figures in this guide and the comparison harnesses drive the editor.

\section{Privacy}
Blendity sends nothing about you or your work anywhere. The one network use it makes on its own is the optional daily
update check. There is no telemetry and no cloud; the only learned model in the program is the optional, local AI
Inbetweening download, which never changes a scene on its own (\code{docs/AI\_DISCLOSURE.md}).

\begin{fromnote}{Learn}
\emph{Welcome to Blendity}: how this editor and the Learn tab work, installing and sharing, updates and older versions,
with or without Blender's libraries.
\end{fromnote}
